\chapter{Laser Skin Treatment}

\section*{Overview}
Laser skin treatment uses focused beams of light (CO2, Erbium, or pulsed dye lasers) to remove damaged outer skin layers or target specific vascular/pigment structures to stimulate collagen growth.

\section*{Alternative Names}
Laser skin resurfacing, Laser dermatology

\section*{Principle}
Laser light of a specific wavelength is absorbed by targeted chromophores in the skin such as pigment, hemoglobin, or water, converting light into heat. The heat destroys or remodels the target tissue while minimizing damage to surrounding skin.
\section*{History}
The ruby laser was first used in dermatology in the 1960s by Leon Goldman. Carbon dioxide (CO2) lasers were developed in the 1980s for skin resurfacing.

\section*{Why It Is Done}
Ordered to treat acne scars, facial wrinkles, sun damage, birthmarks, port-wine stains, or tattoo removal.

\section*{Is This a Primary or Secondary Test or Procedure}
Primary aesthetic and corrective treatment for superficial scarring, vascular lesions, and skin rejuvenation.

\section*{How It Is Performed}
A topical numbing cream is applied. You wear protective goggles. The clinician passes the laser handpiece over the target skin, delivering quick pulses of light.

\section*{Preparation}
Avoid sun exposure and tanning for 4 weeks before the procedure. Stop using retinoids 1 week before.

\section*{Contraindications and Precautions}
Active acne or skin infection in the target area, history of keloid scarring, or active use of isotretinoin (Accutane) within 6 months.

\section*{Risks}
Redness, swelling, hyperpigmentation or hypopigmentation, scarring, or reactivation of herpes simplex (cold sores).

\section*{Aftercare and Recovery}
The skin will feel like a sunburn. Apply gentle moisturizers and strict sun protection. Avoid scratching peeling skin.

\section*{Outcome and Follow-up}
Over weeks, the skin heals, showing smoother texture, reduced pigmentation, and increased elasticity.

\section*{Expected Results}
Successful treatment of the targeted areas; peeling followed by regeneration of healthy, smooth skin.

\section*{Follow-up}
Post-treatment skin check-up; repeat sessions as planned.

\section*{References}
\begin{enumerate}
  \item MedlinePlus. Laser Skin Treatment. U.S. National Library of Medicine; 2025.
  \item Current Medical Diagnosis and Treatment. McGraw-Hill; 2025.
\end{enumerate}

\clearpage
